\documentclass[11pt,a4paper]{article}

\usepackage[T1]{fontenc}
\usepackage[utf8]{inputenc}
\usepackage{lmodern}
\usepackage[margin=2.2cm]{geometry}
\usepackage{array}
\usepackage{booktabs}
\usepackage{longtable}
\usepackage{xcolor}
\usepackage{listings}
\usepackage{hyperref}

\hypersetup{
  colorlinks=true,
  linkcolor=blue!50!black,
  urlcolor=blue!60!black,
  pdftitle={ADDA-CUDA}
}

\lstset{
  basicstyle=\ttfamily\small,
  backgroundcolor=\color{black!3},
  frame=single,
  breaklines=true,
  columns=fullflexible,
  keepspaces=true,
  showstringspaces=false
}

\newcommand{\code}[1]{\nolinkurl{#1}}
\newcommand{\itercmd}[1]{\texttt{-iter #1}}
\newcolumntype{L}[1]{>{\raggedright\arraybackslash}p{#1}}

\title{ADDA-CUDA}
\author{}
\date{}

\begin{document}
\sloppy

\maketitle

CUDA acceleration of \href{https://github.com/adda-team/adda}{ADDA}, the
light-scattering simulator based on the discrete dipole approximation (DDA).

This repository contains a CUDA port developed from the official ADDA source
tree, while preserving the ADDA calculation pipeline and command-line
interface.

ADDA calculates electromagnetic scattering and absorption by particles of
arbitrary shape and composition. It can also handle particles near a plane
substrate and calculate emission or decay-rate enhancement. See the
\href{https://github.com/adda-team/adda}{official ADDA project} and the
\href{https://github.com/adda-team/adda/blob/master/doc/manual.pdf}{ADDA manual} for the physical model and complete list of
options.

\section{Features}

\begin{itemize}
  \item CPU and NVIDIA CUDA implementations of the ADDA solver.
  \item Single precision (\code{float32}) and double precision
        (\code{float64}) executables.
  \item Full 3-D GPU FFT implementation.
  \item Slice-based GPU FFT implementation for reduced VRAM consumption.
  \item Low-memory CUDA organization for larger problems.
  \item CUDA-resident iterative solver vectors, avoiding CPU--GPU transfers
        during the QMR calculation.
  \item CUDA implementations of the available iterative-solver vector
        operations.
  \item Windows build based on Makefiles, MinGW, GCC, G++, and GFortran.
  \item Native Linux build with GCC, FFTW, CMake, and NVIDIA CUDA.
  \item Separate CUDA backend compilation with \code{nvcc}, using the Windows
        \code{.bat} or Linux \code{.sh} script.
  \item Three-level circulant preconditioners developed from Steven Lanier's
        work, including \code{lanier_full}, \code{lanier_partition},
        \code{lanier_nested}, and \code{lanier_multizone}. The \code{ai/}
        directory contains \code{lanier_tqc_v1_ema_actor.bin}, the trained
        actor file created by Steven Lanier and used to configure the
        parameters of the Lanier-preconditioned solver. Reference:
        \href{https://doi.org/10.1016/j.jqsrt.2025.109741}{Steven Lanier,
        ``Learning to precondition: Reinforcement learning enhanced three-level
        circulant preconditioning for the Discrete Dipole Approximation,''
        JQSRT 350 (2026), 109741}. Scientific details are provided in
        \href{adda_cuda_extensions.pdf}{\nolinkurl{adda_cuda_extensions.pdf}},
        \href{adda_cuda_extensions.tex}{\nolinkurl{adda_cuda_extensions.tex}},
        and
        \href{adda_cuda_extensions.md}{\nolinkurl{adda_cuda_extensions.md}}.
  \item Development provenance: the main code was developed from the
        sequential ADDA repository code, without MPI and without OpenCL, with
        assistance from ChatGPT Sol 5.6 in web mode. The CUDA port started from
        upstream ADDA version \code{1.5.0-alpha3}, branch \code{master}, commit
        \href{https://github.com/adda-team/adda/commit/8f550a7786bd4cff5abf3cb2f2180690a08e1157}{\code{8f550a7786bd4cff5abf3cb2f2180690a08e1157}}
        (short SHA \code{8f550a7}), authored by Maxim Yurkin and dated
        30 March 2026 at 08:45:28 UTC, with commit message ``minor changes''. This exact upstream commit is the
        reference baseline for source-code diffs against ADDA-CUDA. The Lanier
        preconditioner code was developed from Steven Lanier's article and from
        written code supplied by Steven Lanier.
\end{itemize}


\section{Split CUDA backend and CLion host debugging}

The default architecture is controlled by:

\begin{lstlisting}
option(ADDA_CUDA_SPLIT_BACKEND
       "Use kernel.cu with the CMake-compiled C++ wrapper backend"
       ON)
\end{lstlisting}

With \code{ADDA_CUDA_SPLIT_BACKEND=ON}, the backend is separated into three
parts. \code{src/kernel.cu} contains only CUDA \code{__global__} kernels,
required \code{__device__} helpers, and C-linkage launch wrappers. It is the
only source compiled by \code{nvcc}. \code{src/kernel.h} defines the private
C ABI used by those launch wrappers. \code{src/wrappermatvec_backend.cpp}
contains the host-side CUDA state, allocations, cuFFT/cuBLAS orchestration,
MatVec dispatch, iterative-solver and Lanier backend logic, and all public
functions declared by \code{src/cudamatvec_backend.h}. CMake compiles this file
with GNU \code{g++}.

The historical \code{src/cudamatvec_backend.cu} is preserved unchanged. When
\code{ADDA_CUDA_SPLIT_BACKEND=OFF}, CMake uses the historical monolithic
backend and the existing libraries under \code{cuda-backend/debug/} or
\code{cuda-backend/release/}.

The public backend ABI is unchanged. The split kernel source is built once for
double precision and once with \code{ADDA_CUDA_SINGLE_BACKEND} for single
precision. Both kernel libraries are stored in \code{cuda-backend/kernels/}.
There is intentionally no separate Debug kernel variant: Debug and Release
CMake profiles share the same nvcc-built kernel library, while the host wrapper
is compiled separately with the normal g++ Debug or Release flags.

On native Windows, nvcc uses a supported MSVC host compiler for
\code{kernel.cu}, while CLion/CMake compile \code{wrappermatvec_backend.cpp}
with MinGW g++. On Linux and WSL, nvcc builds \code{kernel.cu} and the selected
Linux/WSL GNU g++ compiles the wrapper. This allows ordinary source-level
breakpoints and single stepping in the CUDA host code from CLion.

\section{Build scripts and generated files}

With the default split backend, the canonical scripts
\code{scripts/build_cuda_backends.bat} and
\code{scripts/build_cuda_backends.sh} compile only \code{src/kernel.cu}. They
produce one double-precision and one single-precision kernel library in
\code{cuda-backend/kernels/}. The single-precision build defines
\code{ADDA_CUDA_SINGLE_BACKEND}. The same kernel binaries are used by Debug and
Release CMake profiles.

\small
\begin{longtable}{@{}L{0.14\textwidth}L{0.16\textwidth}L{0.31\textwidth}L{0.29\textwidth}@{}}
\toprule
\textbf{Platform} & \textbf{Precision} & \textbf{Kernel runtime} & \textbf{Link input} \\
\midrule
\endhead
Windows & double & \code{cuda-backend/kernels/adda_cuda_kernels.dll} & \code{libadda_cuda_kernels.a} \\
Windows & single & \code{cuda-backend/kernels/adda_cuda_kernels_single.dll} & \code{libadda_cuda_kernels_single.a} \\
Linux/WSL & double & \code{cuda-backend/kernels/libadda_cuda_kernels.so} & same \code{.so} \\
Linux/WSL & single & \code{cuda-backend/kernels/libadda_cuda_kernels_single.so} & same \code{.so} \\
\bottomrule
\end{longtable}
\normalsize

CMake then builds \code{adda_cuda_backend} and
\code{adda_cuda_backend_single} from \code{src/wrappermatvec_backend.cpp} with
GNU g++. The historical Debug/Release backend directories are retained for
\code{ADDA_CUDA_SPLIT_BACKEND=OFF}.

The complete build scripts then create self-contained output directories:

\begin{longtable}{@{}L{0.25\textwidth}L{0.20\textwidth}L{0.47\textwidth}@{}}
\toprule
\textbf{Script} & \textbf{Final directory} & \textbf{Files collected there} \\
\midrule
\endhead
\code{scripts/build_all.bat} & \code{build_windows/bin/} & Eight \code{.exe} files, CUDA DLLs, FFTW DLLs, available MinGW runtime DLLs, and \code{ai/} \\
\code{scripts/build_all.sh} & \code{build_linux/bin/} & Eight Linux executables, both CUDA \code{.so} files, and \code{ai/} \\
\bottomrule
\end{longtable}

\code{build_all.bat} accepts an optional CUDA architecture and MinGW
\code{bin} directory. \code{build_all.sh} accepts an optional CUDA architecture
and number of parallel jobs. When the job count is omitted, each script detects
the number of processors available on the current machine.

\section{Executables}

The complete scripts produce the following executables in
\code{build_windows/bin/} or \code{build_linux/bin/}.

\small
\begin{longtable}{@{}L{0.18\textwidth}L{0.11\textwidth}L{0.14\textwidth}L{0.22\textwidth}L{0.247\textwidth}@{}}
\toprule
\textbf{Executable} & \textbf{Processor} & \textbf{Precision} &
\textbf{FFT / memory organization} & \textbf{Purpose} \\
\midrule
\endhead
\code{adda.exe} & CPU & double / \code{float64} & Standard FFT & Reference ADDA implementation \\
\code{adda_single.exe} & CPU & single / \code{float32} & Standard FFT & CPU single-precision comparison \\
\code{adda_cuda.exe} & NVIDIA GPU & double / \code{float64} & Full 3-D GPU FFT & Accurate CUDA version; highest VRAM use \\
\code{adda_cuda_single.exe} & NVIDIA GPU & single / \code{float32} & Full 3-D GPU FFT & Lower-memory and faster CUDA version \\
\code{adda_cuda_slice.exe} & NVIDIA GPU & double / \code{float64} & Slice FFT & Reduced VRAM in double precision \\
\code{adda_single_slice.exe} & NVIDIA GPU & single / \code{float32} & Slice FFT & Lowest-memory slice variant \\
\code{adda_low_mem.exe} & NVIDIA GPU & single / \code{float32} & Low-memory layout & Minimum large resident GPU buffers \\
\code{adda_low_mem_double.exe} & NVIDIA GPU & double / \code{float64} & Low-memory layout & Low-memory layout with double precision \\
\bottomrule
\end{longtable}
\normalsize

The CUDA backend DLLs are:

\begin{longtable}{@{}L{0.31\textwidth}L{0.17\textwidth}L{0.44\textwidth}@{}}
\toprule
\textbf{DLL} & \textbf{Precision} & \textbf{Role} \\
\midrule
\endhead
\code{adda_cuda_backend.dll} & double / \code{float64} & CUDA backend for double-precision CUDA executables \\
\code{adda_cuda_backend_single.dll} & single / \code{float32} & CUDA backend for single-precision CUDA executables \\
\bottomrule
\end{longtable}

On Linux/WSL with the split backend, the CMake-built host wrapper libraries
are emitted beside the executables and depend on the kernel libraries from
\code{cuda-backend/kernels/}. With \code{ADDA_CUDA_SPLIT_BACKEND=OFF}, the
historical libraries in \code{cuda-backend/debug/} or
\code{cuda-backend/release/} are used instead. The Linux executables have the
same names as in the table above, without the \code{.exe} suffix.

\section{CMake entry points}

The root \code{CMakeLists.txt} contains the shared source lists and build rules.
The platform-specific files only select the appropriate options and call the
root configuration, so the source list is not duplicated:

\begin{lstlisting}
CMakeLists.txt         Shared CMake configuration
windows/CMakeLists.txt Windows / MinGW / CUDA entry point
linux/CMakeLists.txt   Linux / CPU and CUDA entry point
\end{lstlisting}

Use the platform-specific entry point when configuring from a clean checkout.
Do not configure the Windows build from \code{linux/}, or the Linux build from
\code{windows/}.

Complete Windows build with MinGW:

\begin{lstlisting}
scripts\build_all.bat 70 "F:\mingw64_11.2\bin"
\end{lstlisting}

Complete Linux build for \code{sm_70}:

\begin{lstlisting}
bash scripts/build_all.sh 70
\end{lstlisting}

The values \code{70} and \code{sm_70} both select CUDA compute capability 7.0.
It is the default target used in the commands documented above; use the value
matching the target GPU when another architecture is required. If the
architecture argument is omitted, the scripts pass \code{native} to
\code{nvcc} so it can detect the GPU installed on the build machine.

On Windows, \code{"F:/mingw64_11.2/bin"} is the MinGW \code{bin} directory
containing \code{gcc.exe}, \code{g++.exe}, and \code{gfortran.exe}. GFortran is
required by the default complete build. This second argument is optional. When
it is omitted, \code{build_all.bat} finds \code{gcc.exe} from the Windows
\code{PATH} and uses the directory containing it.

The root \code{CMakeLists.txt} remains the shared implementation used by both
entry points.

\section{CLion build}

The project combines three groups of source files:

\begin{itemize}
  \item C and C++ compilation units and include fragments: \code{src/*.c},
        \code{src/*.cpp}, and \code{src/*.inc};
  \item Fortran sources, enabled by default: \code{src/fort/*.f} and
        \code{src/fort/*.f90};
  \item CUDA sources: \code{src/*.cu}.
\end{itemize}

The C sources require the C99 standard, so the Microsoft C/C++ compiler (MSVC)
cannot be used to compile the ADDA C code. On Windows, however, the CUDA
compiler \code{nvcc} uses a supported MSVC toolset as its host compiler. CLion
builds the C, C++, and Fortran sources with the GNU compilers
\code{gcc}, \code{g++}, and \code{gfortran}, using a MinGW, Linux/System, or
WSL toolchain.

For the default split backend, run the kernel build script before configuring
the CUDA targets in CLion:

\begin{lstlisting}[language={}]
scripts\build_cuda_backends.bat
\end{lstlisting}

or, on Linux/WSL:

\begin{lstlisting}[language={}]
bash scripts/build_cuda_backends.sh
\end{lstlisting}

These scripts compile only \code{src/kernel.cu}; no dedicated CUDA Debug
variant is required. CMake compiles \code{src/wrappermatvec_backend.cpp} with
the GNU C++ compiler selected by the CLion toolchain. A Debug profile therefore
provides normal g++ symbols for the CUDA host code, including allocations,
cuFFT/cuBLAS orchestration, MatVec dispatch, iterative solver operations, and
Lanier logic. Kernel source debugging remains a separate CUDA-debugger task.
Set \code{-DADDA_CUDA_SPLIT_BACKEND=OFF} only when intentionally selecting the
preserved historical monolithic backend.

Native Windows builds produce \code{.exe} files. Linux and WSL builds produce
Linux ELF executables, which must be run from Linux or a WSL console. Do not
configure \code{linux/CMakeLists.txt} with MinGW, or
\code{windows/CMakeLists.txt} with Linux/WSL.

\section{Command line build}

\subsection{Linux}

Run the complete script from the repository root. The optional first argument
is the GPU compute capability without the decimal point (\code{70} for
\code{sm_70}, \code{86} for \code{sm_86}, and so on). The optional second
argument is the number of parallel jobs:

\begin{lstlisting}
bash scripts/build_all.sh 70
\end{lstlisting}

This command creates Debug and Release two precision-specific \code{.so} files
in \code{cuda-backend/debug/} and \code{cuda-backend/release/}, configures a
Release build in \code{build_linux/}, builds all eight programs, then copies the shared libraries and the complete
\code{ai/} directory into \code{build_linux/bin/}.

The final directory contains:

\begin{lstlisting}
build_linux/bin/
|-- adda
|-- adda_single
|-- adda_cuda
|-- adda_cuda_single
|-- adda_cuda_slice
|-- adda_single_slice
|-- adda_low_mem
|-- adda_low_mem_double
|-- libadda_cuda_backend.so
|-- libadda_cuda_backend_single.so
`-- ai/
\end{lstlisting}

A short CPU single-precision smoke test is:

\begin{lstlisting}
./build_linux/bin/adda_single -grid 8 8 8 -maxiter 20
\end{lstlisting}

The Linux configuration deliberately uses the system FFTW libraries rather
than the Windows libraries stored in \code{fftw/}. Custom native FFTW paths can
be provided with \code{FFTW3_ROOT}, \code{FFTW3_LIBRARY}, and
\code{FFTW3F_LIBRARY}.

\subsection{Windows}

From a Windows command prompt, run:

\begin{lstlisting}
scripts\build_all.bat 70 "F:\mingw64_11.2\bin"
\end{lstlisting}

The CUDA architecture and MinGW directory are optional. The script first calls
\code{build_cuda_backends.bat}, configures CMake in \code{build_windows/},
builds all eight executables, and collects the runtime files in
\code{build_windows/bin/}.

The resulting directory has the following structure:

\begin{lstlisting}
build_windows/
|-- CMakeCache.txt
|-- CMakeFiles/
|-- Makefile or build.ninja
|-- cmake_install.cmake
`-- bin/
    |-- adda.exe
    |-- adda_single.exe
    |-- adda_cuda.exe
    |-- adda_cuda_single.exe
    |-- adda_cuda_slice.exe
    |-- adda_single_slice.exe
    |-- adda_low_mem.exe
    |-- adda_low_mem_double.exe
    |-- adda_cuda_backend.dll
    |-- adda_cuda_backend_single.dll
    |-- libfftw3-3.dll
    |-- libfftw3f-3.dll
    |-- available MinGW runtime DLLs
    `-- ai/
\end{lstlisting}

The CMake cache, generated build system, and object files remain under
\code{build_windows/}. Only \code{build_windows/bin/} is intended as the
runtime directory. The CUDA import libraries remain in \code{cuda-backend/};
only the CUDA DLLs are copied beside the executables.

To rebuild only the CUDA libraries and their import libraries, run:

\begin{lstlisting}
scripts\build_cuda_backends.bat 70 "F:\mingw64_11.2\bin"
\end{lstlisting}

This backend-only command writes Debug and Release double- and
single-precision DLLs, MSVC \code{.lib} files, and MinGW \code{.a} import
libraries to \code{cuda-backend/debug/} and \code{cuda-backend/release/}. It
does not rebuild the eight ADDA executables.

\section{Requirements}

\subsection{Linux}

\begin{itemize}
  \item C and Fortran compilers supported by the project (GCC and GFortran are
        the tested configuration).
  \item CMake 3.20 or newer and either Make or Ninja.
  \item FFTW development libraries for both precisions. On Debian/Ubuntu,
        install \code{libfftw3-dev}.
  \item For the CUDA targets: an NVIDIA CUDA Toolkit providing \code{nvcc},
        cuFFT, and cuBLAS. A compatible NVIDIA GPU and driver are required to
        run the CUDA executables.
\end{itemize}

For example, on Debian/Ubuntu:

\begin{lstlisting}
sudo apt update
sudo apt install build-essential gfortran cmake libfftw3-dev
\end{lstlisting}

Install the NVIDIA CUDA Toolkit separately if \code{nvcc --version} is not
available. The distribution package name and supported CUDA version depend on
the Linux and NVIDIA driver versions.

\clearpage
\subsection{Windows}

\begin{itemize}
  \item 64-bit Windows.
  \item MinGW-w64 with GCC, GFortran, and \code{dlltool}.
  \item CMake and either MinGW Makefiles or Ninja.
  \item NVIDIA GPU and a compatible NVIDIA driver.
  \item NVIDIA CUDA Toolkit with \code{nvcc} and a supported MSVC host compiler
        (\code{cl.exe}).
  \item The FFTW headers, import libraries, and runtime DLLs stored in
        \code{fftw/}.
\end{itemize}

The root \code{CMakeLists.txt} builds the ADDA C and Fortran sources by default.
\code{scripts/build_all.bat} follows this CMake setting: it does not override
\code{ADDA_NO_FORTRAN}, and puts the selected MinGW \code{bin} directory first
in \code{PATH} so CMake can discover the matching GNU compilers. It also copies
the available GFortran runtime DLLs beside the executables when they are
present. To disable Fortran, change the default in \code{CMakeLists.txt} or
configure with \code{-DADDA_NO_FORTRAN=ON}. The CUDA \code{.cu} source remains
compiled separately by \code{nvcc}, which uses the supported MSVC host compiler
on Windows.

\section{The \texorpdfstring{\texttt{fftw/}}{fftw/} directory}

The \code{fftw/} directory is not included in this GitHub repository. Before
the standard Windows build, the user must download or copy a compatible
Windows FFTW3 distribution into the repository root as \code{fftw/}. CMake
detects this directory automatically when \code{fftw/fftw3.h} is present.
Alternatively, an existing FFTW installation can be selected explicitly with
the CMake variables shown below.

The files most useful for building ADDA on Windows are:

\begin{longtable}{@{}L{0.30\textwidth}L{0.62\textwidth}@{}}
\toprule
\textbf{File} & \textbf{Purpose} \\
\midrule
\endhead
\code{fftw3.h} & FFTW3 C header included by ADDA \\
\code{libfftw3.dll.a} & MinGW import library for double precision \\
\code{libfftw3-3.dll} & Runtime DLL for double-precision CPU/CUDA executables \\
\code{libfftw3f.dll.a} & MinGW import library for single precision \\
\code{libfftw3f-3.dll} & Runtime DLL for single-precision executables \\
\code{libfftw3.a}, \code{libfftw3f.a} & Static libraries, available for alternative link configurations \\
\bottomrule
\end{longtable}

The directory also contains FFTW Fortran interfaces, long-double/quadruple
precision variants, threaded/OpenMP variants, and FFTW wisdom utilities. They
are not required by the default ADDA CMake targets. The CMake build copies the
appropriate \code{libfftw3-3.dll} or \code{libfftw3f-3.dll} next to each
executable.

For a different FFTW installation, override the automatic detection, for
example:

\begin{lstlisting}
cmake -S windows -B build_windows -G Ninja `
  -DFFTW3_ROOT="C:/path/to/fftw" `
  -DFFTW3_LIBRARY="C:/path/to/fftw/libfftw3.dll.a" `
  -DFFTW3F_LIBRARY="C:/path/to/fftw/libfftw3f.dll.a"
\end{lstlisting}

On Linux, the files in this Windows-oriented \code{fftw/} directory should not
be used as native libraries. Install FFTW for Linux or provide Linux-compatible
paths through \code{FFTW3_ROOT}, \code{FFTW3_LIBRARY}, and
\code{FFTW3F_LIBRARY}.

The names above describe the builds tested on Windows on 28 August 2026. The
repository contains source code and build scripts; compiled binaries are not
required to be committed to Git.

\section{Precision and convergence}

Single precision reduces the storage required by complex vectors and FFT
workspaces by approximately a factor of two. It is useful for problems that do
not fit in GPU memory in double precision, but it also reduces the attainable
numerical accuracy and can make iterative convergence more difficult.

Double precision is recommended for validation, difficult or poorly
conditioned problems, and final production results. In single precision,
\code{qmr2} is often preferable to \code{qmr} because it has better round-off
properties. The final residual can be checked with:

\begin{lstlisting}
adda_cuda_single.exe ... -iter qmr2 -recalc_resid
\end{lstlisting}

For CUDA single-precision calculations, the \code{-reliable_resid} option
periodically evaluates the true residual instead of relying only on the
recursively estimated residual. It is intended to limit residual drift and
improve convergence reliability. In the current implementation it applies to
CUDA \code{float32} BiCGStab(2)/BCGS2 and GPBiCGStab(2):

\begin{lstlisting}
adda_cuda_single.exe ... -iter bcgs2 -reliable_resid
adda_cuda_single.exe ... -iter gpbicgstab2 -reliable_resid
\end{lstlisting}

The true residual is checked every 20 iterations. With
\code{-reliable_resid}, the Krylov calculation is restarted only when the
discrepancy indicates a significant residual error. With
\code{-reliable_resid_force_restart}, a restart is forced at every 20-iteration
check, even when the recursively estimated residual is still considered
reliable; this option implies \code{-reliable_resid}.

Examples:

\begin{lstlisting}
adda_cuda_single.exe ... -iter bcgs2 -reliable_resid
adda_cuda_single.exe ... -iter bcgs2 -reliable_resid_force_restart
\end{lstlisting}

These options are ignored for unsupported solvers or for
non-CUDA/non-single-precision builds. They can be combined with
\code{-recalc_resid} when a final independent residual check is required.

\section{Iterative solvers}

Compared with the original ADDA solver set, this version adds iterative
solvers also available in
\href{https://www.fresnel.fr/perso/chaumet/ifdda.html}{IFDDA}:

The command-line syntax is \texttt{-iter SOLVER}; the space between
\texttt{-iter} and the solver name is required.

\begin{longtable}{@{}L{0.31\textwidth}L{0.31\textwidth}L{0.28\textwidth}@{}}
\toprule
\textbf{Solver} & \textbf{Command-line option} & \textbf{Alias} \\
\midrule
\endhead
Bi-CGStab(2) & \itercmd{bcgs2} & \itercmd{bicgstab2} \\
Bi-CG & \itercmd{bicg} & \\
Bi-CGStab & \itercmd{bicgstab} & \\
CGNR & \itercmd{cgnr} & \\
CSYM & \itercmd{csym} & \\
QMR & \itercmd{qmr} & \\
QMR2 & \itercmd{qmr2} & \\
BiCGStab(4) & \itercmd{bicgstab4} & \itercmd{bicgstabl4} \\
BiCGStab(8) & \itercmd{bicgstab8} & \itercmd{bicgstabl8} \\
BiCGStab(12) & \itercmd{bicgstab12} & \itercmd{bicgstabl12} \\
GPBiCGStab(2) & \itercmd{gpbicgstab2} & \itercmd{gpbicgstabl2} \\
GPBiCGStab(4) & \itercmd{gpbicgstab4} & \itercmd{gpbicgstabl4} \\
\bottomrule
\end{longtable}

The default solver remains \code{qmr}. For single-precision calculations,
\code{qmr2} is a useful first choice when round-off limits convergence.

\section{Running ADDA}

ADDA uses command-line options. For example, a 32-cubed grid and refractive
index 1.6 can be requested with options of the following form:

\begin{lstlisting}
adda_cuda_single.exe -grid 32 -m 1.6 0 ...
\end{lstlisting}

The complete syntax, geometry options, material definitions, solver options,
and output files are documented in the \href{https://github.com/adda-team/adda/blob/master/doc/manual.pdf}{ADDA manual}. Use
\code{-h} for the help available in the executable.

\section{Repository layout}

\begin{lstlisting}
src/                       ADDA C sources and CUDA integration
src/kernel.cu             CUDA kernels and nvcc launch wrappers
src/kernel.h              Private C launch-wrapper ABI
src/wrappermatvec_backend.cpp Host CUDA backend compiled by g++
src/cudamatvec_backend.cu Preserved historical monolithic backend
cuda-backend/              Separate CUDA backend build files
scripts/                   Windows and Linux CUDA build scripts
tests/                     Regression and CUDA tests
third_party/               Required third-party headers
doc/                       ADDA documentation
\end{lstlisting}

\section{Relationship to ADDA}

ADDA-CUDA is based on the source code of the official
\href{https://github.com/adda-team/adda}{ADDA project}. The upstream baseline
used to start the CUDA port is ADDA version \code{1.5.0-alpha3}, branch
\code{master}, commit
\href{https://github.com/adda-team/adda/commit/8f550a7786bd4cff5abf3cb2f2180690a08e1157}{\code{8f550a7786bd4cff5abf3cb2f2180690a08e1157}}
(short SHA \code{8f550a7}), authored by Maxim Yurkin and dated 30 March 2026
at 08:45:28 UTC. The commit message is ``minor changes''. Using this exact commit as the upstream baseline
makes it possible to compare ADDA-CUDA with the original ADDA source using a
Git diff. The CUDA implementation is adapted to ADDA's C99 code, FFT pipeline,
complex arithmetic, and iterative solvers.

For the original ADDA scientific references, capabilities, limitations, and
licensing information, please consult the official ADDA repository and the
\href{https://github.com/adda-team/adda/blob/master/doc/manual.pdf}{ADDA manual}.
This project is intended to remain compatible with the upstream ADDA
command-line workflow wherever the CUDA implementation permits.

\section{License and citation}

ADDA is distributed under the GNU General Public License. Please consult the
license and attribution information in the upstream project before
redistributing modified versions.

If results obtained with adda-cuda are published, please cite the original ADDA work and the relevant discrete-dipole-approximation references listed in the official ADDA documentation. Please also reference the adda-cuda repository:
\href{https://github.com/michelgross34/adda-cuda}{ \nolinkurl{https://github.com/michelgross34/adda-cuda}}



For reproducibility, please also report the ADDA-CUDA release tag or commit SHA
used for the published calculations. If a Lanier preconditioner is used, please
also cite Steven Lanier's paper:
\href{https://doi.org/10.1016/j.jqsrt.2025.109741}{``Learning to precondition:
Reinforcement learning enhanced three-level circulant preconditioning for the
Discrete Dipole Approximation,'' JQSRT 350 (2026), 109741}.

\section{Status}

This repository is an experimental CUDA extension of ADDA. CPU and CUDA
results should be compared on representative cases, especially when using
single precision, slice FFT, or low-memory modes. Numerical agreement and
convergence should be validated before using a configuration for production
calculations.

\end{document}
